\chapter{Evaluation}
\label{chap:evaluation}

The generated data are evaluated from two broad perspectives. The first is
\emph{statistical representativeness}: whether the synthetic data reproduce
the marginal distributions, dependence structure, tails, rare categories, and
rare combinations observed in real data. The second is \emph{practical
performance and risk}: whether models trained on synthetic data transfer to
untouched real patients, whether synthetic records remain distinguishable from
real records, and whether they exhibit signs of memorising the training data.
Within these two perspectives, the implementation reports five primary metric
groups: fidelity, tail and rare-event coverage, downstream utility,
real-versus-synthetic detectability, and privacy proxies.

All primary metrics compare a generated dataset with the same frozen real
validation or test partition. The validation partition is used during method
development; the test partition is reserved for the final frozen
configuration. Synthetic datasets contain the same number of rows across
generator variants, and every A1--A5 result is compared with the A0 baseline
under the same split and evaluation protocol.

\section{Statistical Representativeness}
\label{sec:statistical_representativeness}

\subsection{Continuous-Feature Fidelity}
\label{subsec:continuous_fidelity}

Let $X_j^{\mathrm{real}}$ and $X_j^{\mathrm{syn}}$ denote the observed real and
synthetic values of continuous feature $j$. Missing values are removed before
the univariate distances are calculated.

\subsubsection{Wasserstein Distance}

The first Wasserstein distance measures the amount of probability mass that
must be moved to transform one univariate distribution into the other. In
terms of their empirical cumulative distribution functions (ECDFs), it is

\begin{equation}
    W_{1,j} = \int_{-\infty}^{\infty}
        \left|F_j^{\mathrm{real}}(x)-F_j^{\mathrm{syn}}(x)\right|\,\mathrm{d}x.
    \label{eq:evaluation_wasserstein}
\end{equation}

Smaller values indicate better agreement. Because raw Wasserstein distances
depend on the unit and scale of a feature, the headline implementation also
reports an IQR-scaled distance:

\begin{equation}
    W_{1,j}^{\mathrm{scaled}} = \frac{W_{1,j}}{s_j},
    \qquad
    s_j = Q_{.75}(X_j^{\mathrm{real}})-Q_{.25}(X_j^{\mathrm{real}}).
    \label{eq:evaluation_scaled_wasserstein}
\end{equation}

If the real-data IQR is zero, the real-data standard deviation is used; if that
is also zero or non-finite, $s_j$ is set to 1. The mean scaled Wasserstein
distance across continuous features is one of the principal fidelity summary
metrics.

\subsubsection{Kolmogorov--Smirnov Statistic}

The two-sample Kolmogorov--Smirnov statistic measures the largest absolute gap
between the complete real and synthetic ECDFs:

\begin{equation}
    \operatorname{KS}_j = \sup_x
        \left|F_j^{\mathrm{real}}(x)-F_j^{\mathrm{syn}}(x)\right|.
    \label{eq:evaluation_ks}
\end{equation}

It lies in $[0,1]$, with zero indicating identical empirical distributions.
Unlike the Wasserstein distance, it is invariant to the units of the feature.

\subsubsection{Correlation Distance}

To measure whether linear dependence between continuous variables is
preserved, the implementation compares the real and synthetic Pearson
correlation matrices. For $d_c$ continuous features,

\begin{equation}
    D_{\mathrm{corr}} =
        \frac{\left\|R^{\mathrm{real}}-R^{\mathrm{syn}}\right\|_F}{d_c},
    \label{eq:evaluation_correlation}
\end{equation}

where $\|\cdot\|_F$ is the Frobenius norm. Smaller values indicate better
preservation of pairwise correlations.

\subsection{Categorical-Feature Fidelity}
\label{subsec:categorical_fidelity}

Categorical distributions are compared using the Jensen--Shannon distance.
Let $P_j$ and $Q_j$ be the real and synthetic probability vectors over the
union of observed categories, including an explicit missing-value category,
and let $M_j=(P_j+Q_j)/2$. The distance is

\begin{equation}
    \operatorname{JSDist}_j =
    \sqrt{\frac{1}{2}D_{\mathrm{KL}}(P_j\|M_j)
          +\frac{1}{2}D_{\mathrm{KL}}(Q_j\|M_j)}.
    \label{eq:evaluation_js}
\end{equation}

The implementation uses logarithms to base 2, so the result lies in $[0,1]$.
A value of zero means that the category frequencies are identical. The mean
Jensen--Shannon distance across categorical features is reported as the
categorical fidelity summary.

\section{Tail and Rare-Case Representation}
\label{sec:tail_rare_evaluation}

Global distribution distances can conceal clinically important discrepancies
in low-frequency regions. The implementation therefore combines tail-restricted
CDF comparison, quantile mass and location errors, rare-category errors, and
rare-outcome errors.

\subsection{CDF-Tail Divergence}
\label{subsec:cdf_tail_divergence}

For each continuous feature, the main runner evaluates the upper tail beginning
at the real-data 90th percentile. Let
$q_{j,.90}=Q_{.90}(X_j^{\mathrm{real}})$ and
$m_j=\max(X_j^{\mathrm{real}})$. The tail divergence is

\begin{equation}
    D_{j}^{\mathrm{tail}} =
    \max_{x\in\mathcal{G}_j}
    \left|F_j^{\mathrm{real}}(x)-F_j^{\mathrm{syn}}(x)\right|,
    \label{eq:evaluation_cdf_tail}
\end{equation}

where $\mathcal{G}_j$ is a 200-point equally spaced grid over
$[q_{j,.90},m_j]$. If the interval is degenerate, the divergence is set to
zero. Smaller values indicate closer agreement in the upper-tail CDF. The
mean across continuous features is reported as
\texttt{mean\_cdf\_tail\_divergence}.

\subsection{Quantile Mass, Coverage, and Location Error}

The implementation evaluates the lower 5th percentile and the upper 95th and
99th percentiles. For a lower-tail probability $\tau=.05$, define

\begin{equation}
    p_{j,\tau}^{\mathrm{real}}
        =P(X_j^{\mathrm{real}}\le q_{j,\tau}),
    \qquad
    p_{j,\tau}^{\mathrm{syn}}
        =P(X_j^{\mathrm{syn}}\le q_{j,\tau}),
\end{equation}

where $q_{j,\tau}$ is estimated from the real evaluation partition. For upper
quantiles, the probabilities are instead calculated using
$X_j\ge q_{j,\tau}$. The tail-mass error is

\begin{equation}
    E_{j,\tau}^{\mathrm{mass}} =
        \left|p_{j,\tau}^{\mathrm{syn}}-p_{j,\tau}^{\mathrm{real}}\right|.
    \label{eq:evaluation_tail_mass}
\end{equation}

Coverage is the observed synthetic tail mass divided by the nominal expected
tail probability and capped at one:

\begin{equation}
    C_{j,\tau} = \min\left(
        \frac{p_{j,\tau}^{\mathrm{syn}}}{e_{\tau}},1\right),
    \qquad
    e_{\tau}=\begin{cases}
        \tau, & \text{lower tail},\\
        1-\tau, & \text{upper tail}.
    \end{cases}
    \label{eq:evaluation_tail_coverage}
\end{equation}

Thus, coverage measures whether enough synthetic observations reach the real
tail threshold, while the mass error also penalises excessive tail mass. The
scaled quantile-location error is

\begin{equation}
    E_{j,\tau}^{\mathrm{quantile}} =
        \frac{\left|Q_{\tau}(X_j^{\mathrm{syn}})
        -Q_{\tau}(X_j^{\mathrm{real}})\right|}{s_j},
    \label{eq:evaluation_quantile_error}
\end{equation}

where $s_j$ is the same robust scale used in
Equation~\ref{eq:evaluation_scaled_wasserstein}.

\subsection{Rare Categories and Rare Outcomes}

For categorical feature $j$, let

\begin{equation}
    \mathcal{C}^{\mathrm{rare}}_j =
        \left\{c:P_{\mathrm{real}}(X_j=c)\le .05\right\}.
\end{equation}

The rare-category frequency error is

\begin{equation}
    E_j^{\mathrm{rare\mbox{-}cat}} =
        \sum_{c\in\mathcal{C}^{\mathrm{rare}}_j}
        \left|P_{\mathrm{real}}(X_j=c)
        -P_{\mathrm{syn}}(X_j=c)\right|.
    \label{eq:evaluation_rare_category}
\end{equation}

For the prediction target, the code also identifies the least frequent class
in the real evaluation partition and reports its absolute frequency error:

\begin{equation}
    E^{\mathrm{rare\mbox{-}outcome}} =
        \left|P_{\mathrm{real}}(Y=c_{\mathrm{rare}})
        -P_{\mathrm{syn}}(Y=c_{\mathrm{rare}})\right|.
    \label{eq:evaluation_rare_outcome}
\end{equation}

Lower values are preferred for both metrics.

\subsection{RQ1 Fixed-Reference Rare-Region Metrics}
\label{subsec:rq1_rare_regions}

The multi-generator architecture comparison additionally fixes all rare-region
definitions using the real training partition. This avoids allowing each
evaluation sample to move the thresholds. For a continuous feature, the fixed
regions are the values at or below the training 5th percentile and at or above
the training 95th percentile. For a categorical feature, a category is rare
when its training frequency is at most 5\%. At the 5th, 95th, and 99th
continuous-feature quantiles, the code reports the absolute difference between
real-evaluation and synthetic tail mass and the scaled quantile-location error.

A row is also flagged as belonging to a joint rare region when at least two of
the fixed continuous-tail or rare-category conditions hold. If
$Z_{ij}\in\{0,1\}$ indicates that record $i$ occupies rare region $j$, the
joint rare-region mass is

\begin{equation}
    p^{\mathrm{joint}}(D) = \frac{1}{|D|}
        \sum_{i\in D}\mathbb{I}\!\left[\sum_j Z_{ij}\ge 2\right],
    \label{eq:evaluation_joint_rare_mass}
\end{equation}

and the reported error is

\begin{equation}
    E^{\mathrm{joint}} = \left|
        p^{\mathrm{joint}}(D_{\mathrm{eval}})
        -p^{\mathrm{joint}}(D_{\mathrm{syn}})\right|.
    \label{eq:evaluation_joint_rare_error}
\end{equation}

The least frequent target class is determined from the real training
partition. Within this class, the implementation recomputes continuous and
categorical fidelity metrics after excluding the target itself from the
categorical predictors. This outcome-conditional analysis checks whether a
generator represents the feature distributions of rare-outcome patients, not
merely the overall outcome frequency.

Finally, the RQ1 runner estimates a finite-sample reference by drawing two
independent bootstrap samples from the same real evaluation partition and
computing the complete fidelity and rare-region protocol between them. Model
results are reported both directly and relative to this real--real reference.
The reference provides context for sampling variability: even two finite real
samples are not expected to have exactly zero empirical distance.

\subsection{Fixed-Size Rare-Combination Support Coverage}
\label{subsec:support_coverage}

Support coverage is implemented as a supplementary standalone analysis rather
than as an output of the standard A0--A5 evaluation runner. It measures whether
rare multivariate configurations found in real data occur at least once in the
synthetic data without constructing the full joint distribution over all
features.

Numeric variables are discretised using quantile bins learned exclusively from
the real data; categorical variables are matched directly. For a feature
subset $S$ of fixed size $k$, the target label can optionally be appended to
the tuple. Let $\mathcal{R}_S$ be the set of real tuples whose empirical
frequency is at most the configured rarity threshold, and let
$\mathcal{S}_S$ be the synthetic support. Coverage for subset $S$ is

\begin{equation}
    \operatorname{SC}(S) =
        \frac{|\mathcal{R}_S\cap\mathcal{S}_S|}{|\mathcal{R}_S|}.
    \label{eq:evaluation_support_coverage}
\end{equation}

The reported support coverage is the mean over valid $k$-feature subsets; its
standard deviation and the numbers of rare and covered tuples are also saved.
If the number of possible subsets exceeds the configured maximum, subsets are
sampled reproducibly using seed 42. This fixed-size design limits combinatorial
growth and makes coverage more comparable between datasets with different
numbers of columns. Higher coverage is better, although it measures presence
only and does not assess whether the synthetic tuple frequencies are correct.

\section{Practical Performance and Risk}
\label{sec:practical_performance_risk}

\subsection{Downstream Predictive Utility}
\label{subsec:downstream_utility}

The principal utility experiment follows a train-on-synthetic,
test-on-real protocol. A class-balanced random forest is fitted using a
variant's synthetic data and evaluated on the untouched real validation or test
partition. Numeric predictors are median-imputed and robust-scaled;
categorical predictors are most-frequent-imputed and one-hot encoded. The
prediction target is excluded from the predictors. The classifier uses its
fixed argmax decision rule; no decision threshold is selected using audit,
validation, or test labels.

Mortality is evaluated under two training regimes. For the synthetic-only
\texttt{mortality} task, the generated class prevalence is left unchanged.
The \texttt{mortality\_balanced} task resamples the synthetic training data to
a 50/50 class balance. Both regimes are evaluated on the same naturally
imbalanced real partition. The positive label is configured explicitly as 1
rather than being selected manually after inspecting the evaluation data. In
the mixed real/synthetic experiment described below, the unbalanced regime
instead fixes the combined training prevalence to that of the original real
training partition so that mixture fractions are comparable.

Let TP, TN, FP, and FN denote the binary confusion-matrix counts for the
configured positive class. The main threshold-dependent metrics are

\begin{align}
    \operatorname{Accuracy} &=
        \frac{\mathrm{TP}+\mathrm{TN}}
             {\mathrm{TP}+\mathrm{TN}+\mathrm{FP}+\mathrm{FN}},
        \label{eq:evaluation_accuracy}\\
    \operatorname{Precision}_{+} &=
        \frac{\mathrm{TP}}{\mathrm{TP}+\mathrm{FP}},
        \label{eq:evaluation_precision}\\
    \operatorname{Recall}_{+} &=
        \frac{\mathrm{TP}}{\mathrm{TP}+\mathrm{FN}},
        \label{eq:evaluation_recall}\\
    F1_{+} &=
        2\frac{\operatorname{Precision}_{+}\operatorname{Recall}_{+}}
        {\operatorname{Precision}_{+}+\operatorname{Recall}_{+}},
        \label{eq:evaluation_f1}\\
    \operatorname{BalancedAccuracy} &=
        \frac{1}{2}\left(
        \frac{\mathrm{TP}}{\mathrm{TP}+\mathrm{FN}}+
        \frac{\mathrm{TN}}{\mathrm{TN}+\mathrm{FP}}\right).
        \label{eq:evaluation_balanced_accuracy}
\end{align}

Accuracy is retained for completeness, but balanced accuracy, macro F1, and
positive-class recall are more informative for imbalanced mortality data.
Macro precision, recall, and F1 are calculated by treating each class as the
positive class in turn and taking the unweighted mean. This prevents the
majority class from dominating the summary.

The implementation additionally reports ROC-AUC and PR-AUC from predicted
positive-class probabilities. ROC-AUC measures the probability that a randomly
selected positive record receives a higher score than a randomly selected
negative record. PR-AUC summarises the precision--recall curve and is
particularly informative when the positive outcome is rare. For all utility
metrics, larger values indicate better transfer from synthetic training data
to real patients.

The standalone \texttt{RareEventRecall} utility provides a narrower
supplementary analysis. It automatically selects the least frequent real class,
not a manually chosen class, and compares rare-class recall for a real-to-real
random-forest baseline with recall from a synthetic-to-real model. The standard
A0--A5 pipeline supersedes this single-metric view by reporting configured
positive-class recall together with precision, F1, balanced accuracy, ROC-AUC,
and PR-AUC on a common frozen split.

\subsection{Mixed Real/Synthetic Utility}
\label{subsec:mixed_utility_evaluation}

Two mixture protocols test how synthetic data behave when combined with real
training records. Let $N=|D_{\mathrm{train}}|$ and let $f$ be the configured
synthetic fraction parameter.

In the additive protocol, all $N$ real rows are retained and $fN$ synthetic
rows are added. The final training size is $(1+f)N$, with
$f\in\{0,.25,.50,1.0\}$. In the replacement protocol, the total size remains
$N$: $(1-f)N$ real rows are combined with $fN$ synthetic rows, using
$f\in\{0,.25,.50,.75,1.0\}$. The requested class prevalence is enforced at
every fraction. MIMIC-IV uses five deterministic repeats with 300-tree random
forests; WiDS uses three repeats with 100-tree forests to limit runtime. Means
and standard deviations are reported for the same utility metrics defined
above.

\subsection{Real-versus-Synthetic Detectability}
\label{subsec:evaluation_detectability}

For each variant, a balanced probe is formed from equally many real evaluation
records and synthetic records. A class-balanced 300-tree random forest is
trained to predict whether a record is real. Thirty per cent of the probe is
held out using a stratified split. Detector ROC-AUC, average precision, and
accuracy are reported.

For a balanced real-versus-synthetic problem, ROC-AUC $=0.5$ indicates chance
ranking and therefore low detectability. Values approaching 1 indicate that
the distributions are readily separable. Detector AUC should be interpreted by
its distance from 0.5 rather than by applying the usual assumption that a
larger classification AUC is desirable. Average precision and accuracy provide
supporting context but depend more directly on class balance and the fixed 0.5
decision threshold.

\subsection{Memorisation and Privacy Proxies}
\label{subsec:privacy_proxies}

The evaluation includes two families of memorisation-risk indicators. They are
privacy proxies only; they do not establish differential privacy or quantify a
formal privacy budget.

First, the exact-match rate is the fraction of synthetic rows that are
identical across all columns to at least one real training row:

\begin{equation}
    R_{\mathrm{exact}} = \frac{1}{|D_{\mathrm{syn}}|}
        \sum_{\mathbf{x}\in D_{\mathrm{syn}}}
        \mathbb{I}[\mathbf{x}\in D_{\mathrm{train}}].
    \label{eq:evaluation_exact_match}
\end{equation}

Lower values indicate fewer direct copies.

Second, nearest-neighbour distances are calculated after applying a shared
preprocessing transformation fitted on the real training reference. Numeric
variables are median-imputed and robust-scaled, categorical variables are
most-frequent-imputed and one-hot encoded, and Euclidean distances are divided
by the square root of the encoded dimensionality. For query record
$\mathbf{x}$,

\begin{equation}
    d(\mathbf{x},D_{\mathrm{train}}) =
        \frac{1}{\sqrt{p}}
        \min_{\mathbf{z}\in D_{\mathrm{train}}}
        \|\widetilde{\mathbf{x}}-\widetilde{\mathbf{z}}\|_2,
    \label{eq:evaluation_nn_distance}
\end{equation}

where $p$ is the number of transformed columns. The 5th, 50th, and 95th
percentiles are reported separately for synthetic queries and held-out real
queries. Their percentile-wise ratio is

\begin{equation}
    R_q^{\mathrm{NN}} =
        \frac{Q_q(d(D_{\mathrm{syn}},D_{\mathrm{train}}))}
             {Q_q(d(D_{\mathrm{eval}},D_{\mathrm{train}}))},
    \qquad q\in\{.05,.50,.95\}.
    \label{eq:evaluation_nn_ratio}
\end{equation}

A ratio near 1 means that synthetic records are about as close to the training
set as naturally held-out real records. Ratios below 1 indicate unusually close
synthetic records and therefore greater potential memorisation risk. Ratios
above 1 are not automatically preferable, because excessive distance can also
indicate poor fidelity. For WiDS, exact matching remains exhaustive, while
nearest-neighbour work uses deterministic caps of 20,000 reference rows and
10,000 query rows.

\section{Comparison and Reporting}
\label{sec:evaluation_reporting}

For the reweighting experiment, A0 is the primary controlled reference. For a
metric $m$, the raw ablation difference is

\begin{equation}
    \Delta_m^{(a)} = m^{(a)}-m^{(\mathrm{A0})},
    \qquad a\in\{\mathrm{A1},\ldots,\mathrm{A5}\}.
    \label{eq:evaluation_ablation_delta}
\end{equation}

Utility metrics are additionally compared with a real-only model trained on
the original real training partition. The reporting code preserves raw values
and raw deltas, then applies metric-specific direction rules only for
visualisation. For error and distance metrics, a reduction relative to A0 is an
improvement. For utility and coverage metrics, an increase is an improvement.
For detector AUC and average precision, improvement is defined as movement
toward the chance target of 0.5. Privacy-proxy distance ratios are interpreted
relative to the held-out-real reference of 1.

No single metric is sufficient to establish synthetic-data quality. A model
can reproduce marginal distributions while omitting rare combinations, obtain
high utility while remaining easily detectable, or appear private merely
because its samples are unrealistic. Conclusions are therefore based on the
joint pattern across fidelity, tail and rare-case representation, downstream
utility, detectability, and privacy proxies rather than on one aggregate score.
